\subsection{Requisitos Funcionais}

Os requisitos funcionais seguintes decorrem do título do trabalho e
delimitam as capacidades de comunicação a desenvolver.

\begin{description}
    \item[RF01 --- Suporte para I2C:]
    O sistema deve suportar comunicação através de uma interface I2C.

    \item[RF02 --- Suporte para RS232:]
    O sistema deve suportar comunicação através de uma interface RS232.

    \item[RF03 --- Suporte para RS485:]
    O sistema deve suportar comunicação através de uma interface RS485.

    \item[RF04 --- Suporte para CAN:]
    O sistema deve suportar comunicação através de uma interface CAN.

    \item[RF05 --- Suporte para Ethernet:]
    O sistema deve suportar comunicação através de uma interface Ethernet.

    \item[RF06 --- Escrita e transmissão:]
    O sistema deve permitir operações de escrita ou transmissão de dados
    através de cada uma das cinco interfaces previstas, de acordo com as
    características da respetiva tecnologia.

    \item[RF07 --- Leitura e receção:]
    O sistema deve permitir operações de leitura ou receção de dados
    através de cada uma das cinco interfaces previstas, de acordo com as
    características da respetiva tecnologia.
\end{description}

\subsection{Requisitos Não Funcionais}

Os requisitos não funcionais seguintes constituem propostas preliminares,
a validar com os docentes. Os respetivos critérios de aceitação serão
definidos após a consolidação dos requisitos.

\begin{description}
    \item[RNF01 --- Modularidade:]
    A organização do sistema deverá permitir separar as funções associadas
    às diferentes interfaces de comunicação.

    \item[RNF02 --- Facilidade de manutenção:]
    A implementação deverá ser organizada e documentada de forma a
    facilitar a compreensão, a correção e a alteração das suas funções.

    \item[RNF03 --- Extensibilidade:]
    A arquitetura deverá facilitar a introdução futura de novas
    funcionalidades de comunicação, sem exigir a reformulação integral
    do sistema.

    \item[RNF04 --- Fiabilidade:]
    O sistema deverá executar as operações de comunicação de forma
    consistente nas condições de funcionamento que vierem a ser definidas.

    \item[RNF05 --- Facilidade de teste e diagnóstico:]
    A organização do sistema deverá facilitar a verificação das operações
    de comunicação e a identificação de falhas.
\end{description}

\subsection{Restrições do Projeto}

O âmbito conhecido do projeto abrange operações de leitura e escrita
através de I2C, RS232, RS485, CAN e Ethernet. Não estão ainda confirmadas
restrições técnicas específicas relativas à implementação.

A plataforma de execução, a organização concreta das interfaces e as
eventuais funções de interação entre tecnologias permanecem por definir.
Esta ausência de definição corresponde ao estado preliminar do projeto
e não constitui, por si só, uma restrição técnica. As restrições
adicionais deverão ser registadas após confirmação com os docentes.

\subsection{Requisitos a Validar com os Docentes}

As questões seguintes destinam-se a esclarecer o âmbito do sistema e
a estabelecer requisitos verificáveis, sem antecipar decisões de
implementação.

\begin{description}
    \item[QV01 --- Interação entre tecnologias:]
    É necessária conversão ou encaminhamento de dados entre protocolos,
    ou apenas a realização de operações de leitura e escrita através
    de cada interface?

    \item[QV02 --- Protocolos normalizados:]
    É obrigatória a utilização de protocolos normalizados, como Modbus?
    Em caso afirmativo, quais?

    \item[QV03 --- Plataforma de execução:]
    Existe um requisito de utilização de microcontrolador ou é permitida
    uma plataforma com Linux?

    \item[QV04 --- Variante de CAN:]
    O suporte para CAN deve abranger Classical CAN, CAN FD ou ambas
    as variantes?

    \item[QV05 --- Âmbito de Ethernet:]
    O que se entende por suporte para Ethernet neste projeto:
    apenas a camada física ou também comunicação através de TCP,
    UDP ou protocolos de aplicação?

    \item[QV06 --- Placa de circuito impresso:]
    É necessário desenvolver uma placa de circuito impresso
    (PCB) própria?

    \item[QV07 --- Desempenho:]
    Existem métricas obrigatórias de desempenho? Em caso afirmativo,
    quais são os respetivos critérios de aceitação e condições
    de avaliação?
\end{description}
